% =============================================================================
% 05-examples.tex —— 可复制的模板示例汇总（可选）
% -----------------------------------------------------------------------------
% 本文件把所有「写作技法」集中成一组可直接复制的 frame：
%   1. 三色 Block（普通 / 警示 / 示例）
%   2. Block 与列表、双栏的组合
%   3. 定义、定理与证明
%   4. 分栏、图文与表格
%   5. 插图与来源标注
%   6. 公式与分步显示
%   7. 网址、路径与特殊字符
% 不需要这些示例页时，删除 main.tex 中对本文件的 \input 即可，
% 不会影响任何正文内容与目录结构。
%
% 说明：本文件刻意不使用 \subsection，避免生成额外的过渡页；
%       若希望示例也进入目录，把下面的 % ===== 标题改成 \subsection{...} 即可。
% =============================================================================

\section{示例}

% -----------------------------------------------------------------------------
% 1. 三色 Block
% -----------------------------------------------------------------------------
\begin{frame}{Block 样式示例}
  \begin{block}{普通说明：block}
    用于研究问题、定义或方法说明。正文可以包含\textbf{粗体}与公式，例如 $E=mc^2$。
  \end{block}

  \begin{alertblock}{注意事项：alertblock}
    用于假设、限制和需要关注的结论；可用 \alert{强调文字}标出关键条件。
  \end{alertblock}

  \begin{exampleblock}{具体示例：exampleblock}
    用于算例、实验设置或应用场景。例如：对三组数据采用相同指标进行比较。
  \end{exampleblock}
\end{frame}

% -----------------------------------------------------------------------------
% 2. Block 中的列表与双栏
% -----------------------------------------------------------------------------
\begin{frame}{Block 中的列表与双栏}
  \begin{columns}[T,onlytextwidth]
    \column{0.48\textwidth}
    \begin{block}{方法步骤}
      \begin{enumerate}
        \item 明确研究对象。
        \item 统一评价指标。
        \item 对比实验结果。
      \end{enumerate}
    \end{block}

    \column{0.48\textwidth}
    \begin{alertblock}{比较条件}
      \begin{itemize}
        \item 使用相同数据集。
        \item 记录参数设置。
        \item 报告适用范围。
      \end{itemize}
    \end{alertblock}
  \end{columns}

  \begin{exampleblock}{带公式的示例}
    样本均值为 $\bar{x}=\frac{1}{n}\sum_{i=1}^{n}x_i$。
    对样本 $1,2,3$，有 $\bar{x}=2$。
  \end{exampleblock}
\end{frame}

% -----------------------------------------------------------------------------
% 3. 定义、定理与证明
% -----------------------------------------------------------------------------
\begin{frame}{定义、定理与证明}
  \begin{definition}[样本均值]
    对 $n\geq1$ 个实数 $x_1,\ldots,x_n$，定义
    $\bar{x}=\frac{1}{n}\sum_{i=1}^{n}x_i$。
  \end{definition}

  \begin{theorem}[均值界]
    $\min_i x_i\leq\bar{x}\leq\max_i x_i$。
  \end{theorem}

  \begin{proof}
    对每个 $i$，均有 $\min_j x_j\leq x_i\leq\max_j x_j$。
    将不等式求和，再除以 $n$ 即得结论。
  \end{proof}
\end{frame}

% -----------------------------------------------------------------------------
% 4. 分栏与表格
% -----------------------------------------------------------------------------
\begin{frame}{分栏与表格}
  \begin{columns}[T,onlytextwidth]
    \column{0.46\textwidth}
      \textbf{要点}（左栏）
      \begin{enumerate}
        \item 栏宽按内容分配，两栏合计不超过 $1$。
        \item \texttt{[T]} 让两栏顶部对齐；\texttt{onlytextwidth} 避免整体溢出。
        \item 栏内可放列表、公式或图片。
      \end{enumerate}
    \column{0.50\textwidth}
      \small
      \begin{tabularx}{\linewidth}{@{}lY@{}}
        \toprule
        维度 & 说明 \\
        \midrule
        数据 & 填写数据来源与规模 \\
        指标 & 填写评价指标与口径 \\
        结论 & 填写该维度的结论 \\
        \bottomrule
      \end{tabularx}
      \sourceNote{表格使用 \texttt{tabularx} 与 \texttt{Y} 列自适应宽度。}
  \end{columns}
\end{frame}

% -----------------------------------------------------------------------------
% 5. 插图与来源标注
% -----------------------------------------------------------------------------
\begin{frame}{插图与来源标注}
  \begin{columns}[T,onlytextwidth]
    \column{0.54\textwidth}
      \centering
      \slideimage[0.5\textheight]{example-figure.png}
    \column{0.42\textwidth}
      \begin{itemize}
        \item \texttt{\textbackslash slideimage\{文件名\}}：按栏宽等比缩放。
        \item 加可选参数可限制高度，例如 \texttt{[0.45\textbackslash textheight]}。
        \item 图片统一放在 \texttt{figures/}，只写文件名即可。
      \end{itemize}
  \end{columns}
  \sourceNote{来源：\texttt{figures/example-figure.png}（占位图，请替换为实际图片）。}
\end{frame}

% 并排两图：适用于对照实验或前后效果比较
\begin{frame}{并排两图}
  \begin{columns}[T,onlytextwidth]
    \column{0.48\textwidth}
      \centering
      \slideimage[0.42\textheight]{example-figure.png}\\
      {\small 方案 A}
    \column{0.48\textwidth}
      \centering
      \slideimage[0.42\textheight]{example-figure.png}\\
      {\small 方案 B}
  \end{columns}
  \sourceNote{来源：同一占位图重复使用，正式汇报时替换为两组真实结果。}
\end{frame}

% -----------------------------------------------------------------------------
% 6. 公式与分步显示
% -----------------------------------------------------------------------------
\begin{frame}{行内与行间公式}
  \begin{block}{公式写法}
    行内公式写作 $a^2+b^2=c^2$；需要居中并强调时使用行间形式：
    \[
      \bar{x}=\frac{1}{n}\sum_{i=1}^{n}x_i,\qquad
      s^2=\frac{1}{n-1}\sum_{i=1}^{n}(x_i-\bar{x})^2 .
    \]
  \end{block}

  \begin{exampleblock}{常用记号}
    均值 $\bar{x}$、方差 $s^2$、求和 $\sum_{i=1}^{n}$、极限 $\lim_{n\to\infty}$。
  \end{exampleblock}
\end{frame}

\begin{frame}{分步显示}
  \begin{exampleblock}{用 \texttt{\textbackslash pause} 逐条显示}
    \begin{enumerate}
      \item 第一步：先给出问题定义。
      \pause
      \item 第二步：再给出方法。
      \pause
      \item 第三步：最后给出结果与结论。
    \end{enumerate}
  \end{exampleblock}
\end{frame}

% \alert<n> 只在第 n 步高亮，适合逐步引导听众
\begin{frame}{按步骤强调}
  \begin{itemize}
    \item 基线方法：按固定顺序处理输入。
    \item \alert<2>{改进点一}：在第一步加入筛选。
    \item \alert<3>{改进点二}：在第二步加入校准。
    \item \alert<4>{结论}：两项改进共同带来提升。
  \end{itemize}
\end{frame}

% -----------------------------------------------------------------------------
% 7. 网址、路径与特殊字符
% -----------------------------------------------------------------------------
\begin{frame}{网址、路径与特殊字符}
  \begin{itemize}
    \item 网址：\url{https://ctan.org/pkg/beamer}
    \item 路径：\path{figures/example-figure.png}
    \item 行内代码：\texttt{tectonic main.tex}
  \end{itemize}

  \begin{block}{特殊字符需要转义}
    百分号写 \texttt{\textbackslash\%}，与号写 \texttt{\textbackslash\&}，
    下划线写 \texttt{\textbackslash\_}，井号写 \texttt{\textbackslash\#}。
    需要保留原始排版的整段代码，可把 frame 声明为
    \texttt{\textbackslash begin\{frame\}[fragile]} 后使用 verbatim 环境。
  \end{block}
\end{frame}
